%% matrix-content.tex -- shared prose and layout helpers for tests/matrix/*.tex
%% Owned by the CORPUS workstream (implementation contract section 2).
%% tests/test-content.tex is shared and read-only; this is the CORPUS copy.
%%
%% There is no filler here: every sentence is real technical prose about
%% typesetting, page building or numerical methods.
%%
%% ---------------------------------------------------------------------------
%% GEOMETRY CONTRACT OF THE MATRIX DOCUMENTS
%%
%% Every scenario document uses
%%     \usepackage[a4paper,textheight=600pt,textwidth=460pt,top=80pt,
%%                 columnsep=18pt]{geometry}
%% which yields, measured on TeX Live 2025 (build/probe/probe-geo.log):
%%     \textheight   600pt   -> S = \@colht = 600pt on a float-free page
%%     \baselineskip  12pt   -> 50 lines per column
%%     \topskip       10pt   -> first baseline 10pt below the column top
%%     \textwidth    460pt , \columnwidth 221pt , \maxdepth 5pt
%% so the worked numbers of report section 2.8 (S = 600pt, R = 0, baseline
%% 12pt) apply to these documents unchanged.  The page is 597.5pt x 845.0pt,
%% so the type block keeps 68.75pt side margins and ends 165pt above the
%% paper bottom, inside the envelope scripts/check-pdf-bounds.py enforces.
%%
%% \WBMlines{n}{tag} emits exactly n one-line paragraphs, so the height of
%% the material above a request is
%%     h(n) = \topskip + (n-1)\baselineskip = 12n - 2 pt   (plus 0..2.2pt of
%% last-line depth).  Documents that need this to be exact declare
%% \raggedbottom, which removes \flushbottom's interline stretching; the
%% \parskip of article (0pt plus 1pt) is then absorbed by \@makecol's \vfil
%% instead of drifting the measured positions.
%%
%% \WBMbandbox{H}{token}{content} builds a band whose box is exactly H tall
%% with zero depth.  With \widebandrulesoff and \setwidebandsep{12pt} the
%% band footprint of report section 1 is therefore
%%     B = H + 2 * 12pt = H + 24pt
%% exactly, with no dependency on font metrics.
%% ---------------------------------------------------------------------------

\ProvidesFile{matrix-content.tex}[2026/09/03 wideband scenario prose (CORPUS)]

\makeatletter
\newcount\wbm@i
\newcount\wbm@n
\newcount\wbm@j
\newcount\wbm@k
\chardef\wbm@nsent=48

%% ------------------------------------------------------------------ sentences
%% Short declarative sentences (at most 36 characters) so that one sentence
%% plus its token always sets as exactly one line in a 221pt column.
\@namedef{wbm@s@1}{TeX breaks paragraphs into lines.}
\@namedef{wbm@s@2}{The page builder runs on its own.}
\@namedef{wbm@s@3}{Glue carries stretch and shrink.}
\@namedef{wbm@s@4}{A vbox reports height and depth.}
\@namedef{wbm@s@5}{Interline glue uses the last depth.}
\@namedef{wbm@s@6}{The output routine owns box 255.}
\@namedef{wbm@s@7}{Penalties mark legal breakpoints.}
\@namedef{wbm@s@8}{Discardable items go at a break.}
\@namedef{wbm@s@9}{Inserts collect notes and floats.}
\@namedef{wbm@s@10}{Marks travel with split material.}
\@namedef{wbm@s@11}{A column is packed to a set height.}
\@namedef{wbm@s@12}{Badness measures how glue is set.}
\@namedef{wbm@s@13}{The topskip fixes the first line.}
\@namedef{wbm@s@14}{Depth below the last line counts.}
\@namedef{wbm@s@15}{Rules take space but carry no glue.}
\@namedef{wbm@s@16}{Kerns are rigid and never dropped.}
\@namedef{wbm@s@17}{Whatsits record shipout effects.}
\@namedef{wbm@s@18}{Node lists are doubly linked.}
\@namedef{wbm@s@19}{Attributes are inherited by nodes.}
\@namedef{wbm@s@20}{Splitting a box preserves order.}
\@namedef{wbm@s@21}{Newton iteration doubles digits.}
\@namedef{wbm@s@22}{Bisection halves the bracket.}
\@namedef{wbm@s@23}{Quadrature fixes its nodes first.}
\@namedef{wbm@s@24}{Extrapolation cancels error terms.}
\@namedef{wbm@s@25}{A tridiagonal solve is linear.}
\@namedef{wbm@s@26}{Pivoting keeps elimination stable.}
\@namedef{wbm@s@27}{Condition numbers bound the error.}
\@namedef{wbm@s@28}{Floating point addition is inexact.}
\@namedef{wbm@s@29}{Cancellation destroys digits.}
\@namedef{wbm@s@30}{Horner's rule needs n steps.}
\@namedef{wbm@s@31}{A residual is cheaper than error.}
\@namedef{wbm@s@32}{Krylov methods build subspaces.}
\@namedef{wbm@s@33}{Preconditioning reshapes spectra.}
\@namedef{wbm@s@34}{Adaptive steps track local error.}
\@namedef{wbm@s@35}{Stiff systems need implicit steps.}
\@namedef{wbm@s@36}{Explicit stages reuse old slopes.}
\@namedef{wbm@s@37}{A grid function samples a domain.}
\@namedef{wbm@s@38}{Discretisation error scales with h.}
\@namedef{wbm@s@39}{Convergence order needs refinement.}
\@namedef{wbm@s@40}{Stability limits the time step.}
\@namedef{wbm@s@41}{Two columns share one page list.}
\@namedef{wbm@s@42}{Balancing redistributes lines.}
\@namedef{wbm@s@43}{A strip spans the full text width.}
\@namedef{wbm@s@44}{Reading order follows content order.}
\@namedef{wbm@s@45}{Floats are deferred until they fit.}
\@namedef{wbm@s@46}{The gutter separates two columns.}
\@namedef{wbm@s@47}{Hyphenation adds optional breaks.}
\@namedef{wbm@s@48}{A ledger proves nothing is lost.}

%% \wbm@sent{<count register value>}: sentence ((k-1) mod 48) + 1
\def\wbm@sent#1{%
  \wbm@k=#1\relax
  \advance\wbm@k\m@ne
  \wbm@j=\wbm@k
  \divide\wbm@j\wbm@nsent
  \multiply\wbm@j\wbm@nsent
  \advance\wbm@k-\wbm@j
  \advance\wbm@k\@ne
  \csname wbm@s@\number\wbm@k\endcsname}

%% \wbm@two{<number>}: two-digit, so no token is a prefix of another
\def\wbm@two#1{\ifnum#1<10 0\fi\number#1}

%% ------------------------------------------------------------ exact-height prose
%% \WBMlines{<n>}{<tag>}: exactly <n> one-line paragraphs.  Line i starts with
%% the unique token <tag><ii>, so h(n) = 12n - 2 pt as documented above.
\newcommand\WBMlines[2]{%
  \par
  \begingroup
    \parindent\z@
    \wbm@n=#1\relax
    \wbm@i=\z@
    \loop\ifnum\wbm@i<\wbm@n
      \advance\wbm@i\@ne
      \noindent\mbox{#2\wbm@two{\wbm@i}}\space\wbm@sent{\wbm@i}\par
    \repeat
  \endgroup}

%% ------------------------------------------------------------- flowing prose
%% \WBMpara{<n>}{<tag>}: <n> ordinary paragraphs of six sentences each, for the
%% scenarios where realistic line breaking matters more than an exact height.
\newcommand\WBMpara[2]{%
  \par
  \begingroup
    \wbm@n=#1\relax
    \wbm@i=\z@
    \loop\ifnum\wbm@i<\wbm@n
      \advance\wbm@i\@ne
      \mbox{#2\wbm@two{\wbm@i}}\space
      \wbm@k=\wbm@i \multiply\wbm@k6 \advance\wbm@k-5
      \wbm@sent{\wbm@k}\space
      \advance\wbm@k\@ne \wbm@sent{\wbm@k}\space
      \advance\wbm@k\@ne \wbm@sent{\wbm@k}\space
      \advance\wbm@k\@ne \wbm@sent{\wbm@k}\space
      \advance\wbm@k\@ne \wbm@sent{\wbm@k}\space
      \advance\wbm@k\@ne \wbm@sent{\wbm@k}%
      \par
    \repeat
  \endgroup}

%% ------------------------------------------------------------------ the bands
%% \WBMbandbox{<H>}{<token>}{<visible content>}: a band body of height exactly
%% <H> and depth 0.  The \vrule strut sets the height; \raisebox with zero
%% declared height and depth centres the visible content without contributing
%% to the measurement.
\newcommand\WBMbandbox[3]{%
  \noindent\hbox to\textwidth{%
    \vrule width\z@ height#1 depth\z@
    \hfil
    \raisebox{\dimexpr#1/2-5pt\relax}[\z@][\z@]{#3\quad\mbox{#2}}%
    \hfil}}

%% Realistic band bodies whose height comes from the mathematics.  Their
%% natural heights on TeX Live 2025 are noted for the derivations; documents
%% that need an exact B use \WBMbandbox instead.
%% \WBMbandeqA  ~ 29pt   \WBMbandeqB  ~ 59pt   \WBMbandeqC  ~ 104pt
\newcommand\WBMbandeqA[1]{%
  \begin{gather*}
    \mathcal{T}[u] = \int_{\Omega} \bigl( \tfrac12 |\nabla u|^2 - f u \bigr) \, dx
  \end{gather*}
  \noindent\hbox to\textwidth{\hfil\mbox{#1}\hfil}}
\newcommand\WBMbandeqB[1]{%
  \begin{gather*}
    \mathcal{T}[u] = \int_{\Omega} \bigl( \tfrac12 |\nabla u|^2 - f u \bigr) \, dx \\
    \delta\mathcal{T}[u;v] = \int_{\Omega} \bigl( \nabla u \cdot \nabla v - f v \bigr) \, dx \\
    A_{ij} = \int_{\Omega} \nabla \varphi_i \cdot \nabla \varphi_j \, dx
  \end{gather*}
  \noindent\hbox to\textwidth{\hfil\mbox{#1}\hfil}}
\newcommand\WBMbandeqC[1]{%
  \begin{gather*}
    \mathcal{T}[u] = \int_{\Omega} \bigl( \tfrac12 |\nabla u|^2 - f u \bigr) \, dx \\
    \delta\mathcal{T}[u;v] = \int_{\Omega} \bigl( \nabla u \cdot \nabla v - f v \bigr) \, dx \\
    A_{ij} = \int_{\Omega} \nabla \varphi_i \cdot \nabla \varphi_j \, dx \\
    b_i = \int_{\Omega} f \varphi_i \, dx \\
    \kappa(A) = \|A\| \, \|A^{-1}\| \\
    \|u - u_h\|_{1} \le C h^{p} \, \|u\|_{p+1}
  \end{gather*}
  \noindent\hbox to\textwidth{\hfil\mbox{#1}\hfil}}

%% ------------------------------------------------------- unbreakable material
%% \WBMblock{<H>}{<token>}: a block of height <H> with no legal breakpoint
%% inside it, used for the D3 "no legal cut" scenario.  A single \vbox is a
%% single item on the vertical list, so the page builder cannot cut into it.
\newcommand\WBMblock[2]{%
  \par\nobreak
  \vbox to#1{%
    \hsize\columnwidth
    \vss
    \noindent\hbox to\columnwidth{\hfil\mbox{#2}\hfil}%
    \vss}%
  \par\nobreak}

%% \WBMnote{<token>}: a footnote whose text carries a unique token.
\newcommand\WBMnote[1]{\footnote{\mbox{#1} Inserts collect notes and floats.}}

%% \WBMlongnote{<n>}{<token>}: a footnote of <n> sentences, tall enough to be
%% split at a column boundary so that a hold-over becomes pending.
\newcommand\WBMlongnote[2]{%
  \footnote{\mbox{#2}\space
    \wbm@n=#1\relax
    \wbm@i=\z@
    \loop\ifnum\wbm@i<\wbm@n
      \advance\wbm@i\@ne
      \wbm@sent{\wbm@i}\space
    \repeat}}

%% \WBMmargin{<token>}: a marginal note narrow enough for the margin geometry
%% the marginpar scenarios declare.
\newcommand\WBMmargin[1]{\marginpar{\tiny\raggedright\mbox{#1}}}

%% \WBMfig{<placement>}{<token>}: a single-column float of known height.
\newcommand\WBMfig[2]{%
  \begin{figure}[#1]
    \centering
    \rule{.7\columnwidth}{40pt}
    \caption{\mbox{#2} A single-column diagnostic float.}
  \end{figure}}

%% \WBMwidefig{<placement>}{<token>}: a double-column float.
\newcommand\WBMwidefig[2]{%
  \begin{figure*}[#1]
    \centering
    \rule{.6\textwidth}{40pt}
    \caption{\mbox{#2} A double-column diagnostic float.}
  \end{figure*}}
\makeatother
\endinput
